\section{Further questions}\label{sec:questions}
%% ====================================================================

The preceding sections complete the divisor-lattice theorem chain for this
article.  We record only two boundary markers; neither is used as evidence for
the results proved above.
\begin{enumerate}[label=(\arabic*)]
\item Classify connected support kernels of size at least four. By the
\suppref{Corollary on four-support reduction}, the one-, two-, and
three-coordinate classifications already determine every disconnected
four-coordinate pattern. The \suppref{Corollary ``Theta-refined connected
support entropy''} shows that connected kernels already realize its lower bound in
the fixed-choice rank-pure minimal-cover sector, so a full connected
high-support classification is not a lower-order residual problem within that
sector. It would be a separate project.
\item Prove or disprove either rank-window sparsity conjecture. For the mean
hypothesis, Theorem~\ref{thm:mean-rank-window-deaggregation} reduces the
problem to the visible single-rank maximum
\eqref{eq:visible-single-rank-interface}; Proposition
\ref{prop:blms-obstructs-mean-window} also shows that a published conjecture
predicts failure rather than sparsity. As explained in
Remark~\ref{rem:exact-average-open-interface}, the remaining question is an
upper distribution problem for exact-rank prime multiplicities.
The paper derives conditional normal and mean orders, but does not
claim that either conjecture follows from presently known primitive-divisor
theorems, RH, or ERH.
\end{enumerate}
